\chapter{How to use this guide}
\label{ch:map}

This note teaches the mathematics that Homeworks 1 and 2 actually use. Each chapter states a definition, gives the intuition, derives the identity, draws the picture, and then poses a short practice problem with a worked solution. The practice problems are smaller than the homework. They train the same moves.

\begin{intbox}[What to read for which homework]
Homework 1 uses Chapters~\ref{ch:prob}--\ref{ch:ar}. Homework 2 uses those chapters and then Chapters~\ref{ch:vae}--\ref{ch:flow}. The conditional ELBO, the hierarchical ELBO, dequantization, and the Glow convolution are extensions of the lectures. Their derivations are written out here.
\end{intbox}

\begin{center}
\renewcommand{\arraystretch}{1.25}
\begin{tabularx}{\textwidth}{@{}lXl@{}}
\toprule
\textbf{Homework piece} & \textbf{Read} & \textbf{Slides}\\
\midrule
Gaussians, products, conditionals & Ch.~\ref{ch:gauss} & L2, 12--35\\
KL, Jensen--Shannon, perplexity & Ch.~\ref{ch:info} & L2, 36--42\\
Monte Carlo, importance sampling & Ch.~\ref{ch:mc} & L2, 34--35; L4, 23--29\\
Markov models and autoregression & Ch.~\ref{ch:ar} & L3\\
ELBO, reparameterization, CVAE & Ch.~\ref{ch:vae} & L4, 12--40, 56--84\\
Two-layer hierarchical ELBO & Ch.~\ref{ch:hier} & pattern of L4, 38--39\\
Coupling, flow prior, dequantization, Glow & Ch.~\ref{ch:flow} & L5, 3--73\\
\bottomrule
\end{tabularx}
\end{center}

\begin{defbox}[Notation used throughout]
Random variables are capitals when the distinction matters, and the same letter is used for a value when the slides do so. A Gaussian $\N(x\given\mu,\sigma^2)$ has variance $\sigma^2$, matching Lecture~2. $\log$ is the natural logarithm, so information quantities are in nats. $\KL(q\Vert p)$ means we use $q$ where the reference in the denominator is $p$. The Jacobian $J_f(u)$ is the matrix of partial derivatives $\partial f_i/\partial u_j$.
\end{defbox}
